% =====================================================================
%  Listas y arreglos en Python — de la lista nativa a NumPy
%  Programación orientada a la computación científica
%  Escuela de Ingeniería de Sistemas e Informática — UIS
%
%  Compilación:
%     pdflatex listas_numpy.tex     (dos veces)
%
%  Imágenes: las mismas de la versión Marpit, en ./img
%     Por omisión se usan los PDF generados con ./svg2pdf.sh
%     Para usar los SVG directamente, ponga \usarsvgtrue abajo
%     y compile con:  pdflatex -shell-escape listas_numpy.tex
%     (requiere Inkscape instalado)
%
%  Notas del presentador: descomente la línea \setbeameroption de abajo,
%     o use pdfpc con las notas embebidas.
% =====================================================================

\PassOptionsToPackage{table}{xcolor}
\documentclass[aspectratio=169, 11pt]{beamer}

% ---------- ¿SVG directo o PDF preconvertido? ----------
\newif\ifusarsvg
\usarsvgfalse            % <- \usarsvgtrue para usar los .svg con Inkscape

\usepackage[T1]{fontenc}
\usepackage[utf8]{inputenc}
\usepackage[spanish,es-noshorthands]{babel}   % requiere texlive-lang-spanish
\usepackage{helvet}
\renewcommand{\familydefault}{\sfdefault}
\usepackage{amsmath}
\usepackage{graphicx}
\usepackage{booktabs}
\usepackage{tabularx}
\usepackage{listings}
\usepackage{tcolorbox}
\tcbuselibrary{skins}
\usepackage{microtype}

\ifusarsvg
  \usepackage{svg}
  \svgpath{{img/}}
\else
  \graphicspath{{img/}}
\fi

% \setbeameroption{show notes on second screen=right}

% ---------------------------------------------------------------------
%  Paleta borgoña
% ---------------------------------------------------------------------
\definecolor{borgonaOscuro}{HTML}{590D22}
\definecolor{borgonaMedio}{HTML}{7B2D3B}
\definecolor{borgonaClaro}{HTML}{A4133C}
\definecolor{borgonaPalido}{HTML}{F2C4CE}
\definecolor{fondoClaro}{HTML}{F9F4F5}
\definecolor{textoOscuro}{HTML}{1A1A1A}
\definecolor{textoMedio}{HTML}{3D3D3D}
\definecolor{lineaSuave}{HTML}{E3CBD1}

% ---------------------------------------------------------------------
%  Apariencia general
% ---------------------------------------------------------------------
\usetheme{default}
\useinnertheme{default}
\setbeamertemplate{navigation symbols}{}
\setbeamercolor{normal text}{fg=textoOscuro, bg=white}
\setbeamercolor{structure}{fg=borgonaOscuro}
\setbeamercolor{alerted text}{fg=borgonaClaro}
\setbeamerfont{frametitle}{size=\Large, series=\bfseries}

% Título de diapositiva: borgoña con filete inferior
\setbeamertemplate{frametitle}{%
  \vspace*{0.6em}%
  {\color{borgonaOscuro}\usebeamerfont{frametitle}\insertframetitle\par}%
  \vspace*{0.25em}%
  {\color{borgonaOscuro}\rule{\textwidth}{2.2pt}}%
  \vspace*{0.2em}%
}

% Paginación
\setbeamertemplate{footline}{%
  \hfill{\color{borgonaMedio}\footnotesize\insertframenumber\,/\,\inserttotalframenumber}%
  \hspace*{2.2em}\vspace*{0.55em}%
}

% Viñetas
\setbeamercolor{itemize item}{fg=borgonaClaro}
\setbeamercolor{itemize subitem}{fg=borgonaMedio}
\setbeamercolor{enumerate item}{fg=borgonaClaro}
\setbeamertemplate{itemize item}{\textbullet}
\setbeamertemplate{itemize subitem}{--}
\setbeamertemplate{enumerate item}{\insertenumlabel.}
\setlength{\leftmargini}{1.4em}

% Código en línea (evita el relleno de \lstinline con columns=fixed)
\newcommand{\code}[1]{{\ttfamily\small\color{textoOscuro}#1}}

% Subtítulos internos (equivalen a h2 del Marpit)
\newcommand{\subt}[1]{%
  \par\smallskip
  \noindent{\color{borgonaClaro}\rule{3pt}{0.95em}}\hspace{0.5em}%
  {\color{borgonaMedio}\bfseries #1}\par\vspace{0.3em}
}

% Cita destacada (equivale a blockquote)
\newtcolorbox{cita}{
  colback=fondoClaro, colframe=borgonaClaro,
  boxrule=0pt, leftrule=3.5pt,
  left=9pt, right=8pt, top=5pt, bottom=5pt,
  sharp corners, fontupper=\itshape\small\color{textoMedio}
}

% ---------------------------------------------------------------------
%  Código
% ---------------------------------------------------------------------
\lstdefinestyle{pyborgona}{
  language=Python,
  basicstyle=\ttfamily\footnotesize\color{textoOscuro},
  keywordstyle=\color{borgonaClaro}\bfseries,
  commentstyle=\color{borgonaMedio}\itshape,
  stringstyle=\color{borgonaOscuro},
  numberstyle=\tiny\color{textoMedio},
  backgroundcolor=\color{fondoClaro},
  frame=l,
  framerule=3pt,
  rulecolor=\color{borgonaMedio},
  framesep=8pt,
  xleftmargin=10pt,
  aboveskip=8pt, belowskip=6pt,
  columns=fixed, basewidth=0.55em,
  keepspaces=true,
  showstringspaces=false,
  breaklines=true,
  extendedchars=true,
  literate=%
    {á}{{\'a}}1 {é}{{\'e}}1 {í}{{\'i}}1 {ó}{{\'o}}1 {ú}{{\'u}}1
    {Á}{{\'A}}1 {É}{{\'E}}1 {Í}{{\'I}}1 {Ó}{{\'O}}1 {Ú}{{\'U}}1
    {ñ}{{\~n}}1 {Ñ}{{\~N}}1 {ü}{{\"u}}1 {¿}{{?`}}1 {¡}{{!`}}1
    {°}{{$^\circ$}}1 {²}{{$^2$}}1 {½}{{$\frac{1}{2}$}}1
    {✓}{{\color{borgonaMedio}$\checkmark$}}1
    {✗}{{\color{borgonaClaro}$\times$}}1
    {→}{{$\rightarrow$}}1 {≥}{{$\geq$}}1 {≈}{{$\approx$}}1
    {—}{{---}}1 {π}{{$\pi$}}1 {σ}{{$\sigma$}}1 {…}{{\ldots}}3
}
\lstset{style=pyborgona}

% ---------------------------------------------------------------------
%  Figuras: puente entre SVG y PDF preconvertido
% ---------------------------------------------------------------------
% #1 = fracción de la altura útil de la diapositiva (por omisión 0,64)
\newcommand{\figura}[2][0.64]{%
  \begin{center}
  \ifusarsvg
    \includesvg[height=#1\textheight]{#2}%
  \else
    \includegraphics[width=\textwidth, height=#1\textheight, keepaspectratio]{#2}%
  \fi
  \end{center}
}

% ---------------------------------------------------------------------
%  Portada y separadores de sección
% ---------------------------------------------------------------------
\newcommand{\portada}{%
{%
\setbeamercolor{background canvas}{bg=borgonaOscuro}
\begin{frame}[plain]
  \vfill
  {\color{white}\fontsize{26}{30}\selectfont\bfseries Listas y arreglos en Python\par}
  \vspace{0.4em}
  {\color{white}\rule{\textwidth}{2.5pt}}
  \vspace{0.9em}

  {\color{borgonaPalido}\Large De la lista nativa a NumPy, para computación científica\par}
  \vspace{2.2em}

  {\color{borgonaPalido}\small
   Sergio Augusto Gélvez Cortés\\[2pt]
   Escuela de Ingeniería de Sistemas e Informática\\[2pt]
   Universidad Industrial de Santander\\[10pt]
   Programación orientada a la computación científica\par}
  \vfill
\end{frame}
}%
\note{Sesión dirigida a estudiantes de física. La meta no es memorizar métodos,
sino que entiendan dos cosas: que una lista es un contenedor de referencias, y
que un ndarray es un bloque de memoria tipado. Casi todo lo demás se deduce de ahí.}
}

\newcommand{\seccion}[2]{%
{%
\setbeamercolor{background canvas}{bg=borgonaMedio}
\begin{frame}[plain]
  \vfill
  {\color{white}\fontsize{24}{28}\selectfont\bfseries #1\par}
  \vspace{0.4em}
  {\color{borgonaPalido}\rule{\textwidth}{2.5pt}}
  \vspace{0.8em}

  {\color{borgonaPalido}\Large #2\par}
  \vfill
\end{frame}
}%
}

% =====================================================================
\begin{document}

\portada

% ---------------------------------------------------------------------
\begin{frame}{Mapa de la sesión}
\begin{enumerate}
  \item \textbf{La lista de Python} — qué es, cómo se crea, cómo se recorre
  \item \textbf{Operaciones básicas} — consultar y modificar
  \item \textbf{Alias y copias} — el error más costoso del semestre
  \item \textbf{Rebanado} — la sintaxis que después reutilizaremos en NumPy
  \item \textbf{Listas anidadas} — matrices con listas, y por qué incomodan
  \item \textbf{NumPy y el ndarray} — qué cambia en la memoria
  \item \textbf{Vectorización, broadcasting y reducciones}
  \item \textbf{Comparación y criterio} — cuándo usar cada uno
\end{enumerate}
\end{frame}
\note{Anunciar desde el principio que la sección de listas no es un rodeo: la
sintaxis de rebanado y de recorrido se hereda casi íntegra en NumPy, y hay
tareas donde la lista sigue siendo la herramienta correcta.}

% =====================================================================
\seccion{1. La lista de Python}{Una secuencia ordenada y mutable}

% ---------------------------------------------------------------------
\begin{frame}[fragile]{¿Qué es una lista?}
Una \textbf{lista} es una secuencia \textbf{ordenada} de referencias a objetos.
Tres propiedades la definen:
\begin{itemize}
  \item \textbf{Ordenada}: cada elemento tiene una posición fija, accesible por índice. No es un conjunto.
  \item \textbf{Mutable}: se puede cambiar su contenido y su longitud después de creada.
  \item \textbf{Heterogénea}: puede contener objetos de tipos distintos, incluso otras listas.
\end{itemize}

\begin{lstlisting}
medidas   = [20.1, 20.4, 19.8, 20.0]     # homogénea, por disciplina
mezclada  = [1, "electrón", 3.14, True]  # legal, pero rara vez deseable
vacia     = []
\end{lstlisting}
\end{frame}
\note{La heterogeneidad es legal pero casi siempre indeseable en cálculo
numérico: si una lista de flotantes tiene un None escondido, el error aparece
a mitad del procesamiento y es difícil de rastrear.}

% ---------------------------------------------------------------------
\begin{frame}{Creación e indexación}
\figura[0.60]{indexacion}
\end{frame}
\note{Insistir en que el índice negativo no es una excentricidad: v[-1] es la
forma idiomática de pedir el último elemento, sin calcular len(v)-1.}

% ---------------------------------------------------------------------
\begin{frame}[fragile]{Operaciones que consultan}
No modifican la lista; devuelven un valor nuevo.

\begin{lstlisting}
v = [3.2, 1.7, 4.8, 0.5, 2.9]

len(v)          # 5      cantidad de elementos
v[2]            # 4.8    acceso por índice
4.8 in v        # True   pertenencia
v.index(4.8)    # 2      posición de la primera aparición
v.count(1.7)    # 1      cuántas veces aparece
min(v), max(v)  # (0.5, 4.8)
sum(v)          # 13.1
sorted(v)       # [0.5, 1.7, 2.9, 3.2, 4.8]   lista NUEVA, v intacta
\end{lstlisting}
\end{frame}
\note{sum, min y max son funciones incorporadas del lenguaje, no métodos de la
lista. sorted devuelve una lista nueva; v.sort() ordena en el sitio.
Confundirlas es un error clásico.}

% ---------------------------------------------------------------------
\begin{frame}[fragile]{Operaciones que modifican}
Actúan \textbf{sobre la lista misma} y devuelven \code{None}.

\begin{center}
\rowcolors{2}{fondoClaro}{white}
\small
\begin{tabularx}{0.94\textwidth}{lX}
\rowcolor{borgonaOscuro}
\color{white}\textbf{Operación} & \multicolumn{1}{l}{\color{white}\textbf{Efecto}} \\
\code{v.append(x)}          & añade \code{x} al final \\
\code{v.extend(otra)}       & añade todos los elementos de otra secuencia \\
\code{v.insert(i, x)}       & inserta \code{x} en la posición \code{i} \\
\code{v.pop()} / \code{v.pop(i)} & extrae y \textbf{devuelve} el último, o el de índice \code{i} \\
\code{v.remove(x)}          & elimina la primera aparición de \code{x} \\
\code{del v[i]}             & elimina el elemento de índice \code{i} \\
\code{v.sort()}             & ordena en el sitio \\
\code{v.reverse()}          & invierte en el sitio \\
\code{v.clear()}            & deja la lista vacía \\
\end{tabularx}
\end{center}
\end{frame}
\note{pop es la única de la tabla que devuelve algo útil. El resto devuelve
None, y ese es justamente el origen del error de la siguiente diapositiva.}

% ---------------------------------------------------------------------
\begin{frame}[fragile]{El error que todos cometen una vez}
\begin{lstlisting}
v = [3, 1, 2]

v = v.sort()      # ✗  v ahora vale None
print(v)          #    None

v = [3, 1, 2]
v.sort()          # ✓  v vale [1, 2, 3]
print(v)          #    [1, 2, 3]

w = sorted([3, 1, 2])   # ✓  también válido: w es una lista nueva
\end{lstlisting}

\begin{cita}
Los métodos que modifican en el sitio devuelven \texttt{None}. Asignar su
resultado destruye la lista.
\end{cita}

Lo mismo aplica a \code{append}, \code{extend}, \code{insert},
\code{reverse} y \code{clear}.
\end{frame}
\note{Vale la pena escribirlo en el tablero y dejarlo un minuto. Es el error
que más tiempo hace perder en los primeros talleres, y el mensaje de Python
("NoneType object is not subscriptable") no señala la causa.}

% ---------------------------------------------------------------------
\begin{frame}[fragile]{Concatenar y repetir}
\begin{lstlisting}
a = [1, 2]
b = [3, 4]

a + b        # [1, 2, 3, 4]       lista nueva
a * 3        # [1, 2, 1, 2, 1, 2] lista nueva
a += b       # equivale a a.extend(b): modifica a en el sitio
\end{lstlisting}

\textbf{Aviso para más adelante:} en una lista, \code{+} concatena y
\code{*} repite. En un arreglo de NumPy, \code{+} suma elemento a
elemento y \code{*} multiplica elemento a elemento. Es la misma sintaxis
con dos significados distintos.
\end{frame}
\note{Este contraste conviene sembrarlo ya, porque es la confusión número uno
al pasar de listas a NumPy: [1,2] + [3,4] da [1,2,3,4], pero
np.array([1,2]) + np.array([3,4]) da [4,6].}

% ---------------------------------------------------------------------
\begin{frame}{Asignar no copia}
\figura[0.66]{alias-copia}
\end{frame}
\note{Aquí conviene mostrar id(a) e id(b) en vivo: son el mismo número en el
caso del alias, distintos en el de la copia. Y mencionar que este
comportamiento es idéntico cuando se pasa una lista a una función: la función
recibe el mismo objeto y puede modificarlo.}

% ---------------------------------------------------------------------
\begin{frame}[fragile]{Alias, copia superficial y copia profunda}
\begin{lstlisting}
import copy

a = [[1, 2], [3, 4]]

b = a                   # alias: mismo objeto
c = a.copy()            # copia superficial (igual que a[:] o list(a))
d = copy.deepcopy(a)    # copia profunda: replica las listas internas

c[0][0] = 99            # ¡también cambia a! la lista interna es compartida
d[0][0] = 77            # a no se entera
\end{lstlisting}

Para el curso: si va a modificar una lista dentro de una función y no quiere
efectos colaterales, \textbf{copie explícitamente al entrar}.
\end{frame}
\note{La copia superficial copia el arreglo de referencias, no los objetos
apuntados. Con listas de números no se nota, porque los números son
inmutables; con listas de listas, sí. Es exactamente el mismo fenómeno que
veremos en NumPy con vistas y copias.}

% ---------------------------------------------------------------------
\begin{frame}{Rebanado (\textit{slicing})}
\figura[0.70]{slicing}
\end{frame}
\note{Enfatizar la asimetría: inicio incluido, fin excluido. Es la misma
convención de range() y garantiza que v[:k] + v[k:] reconstruye v para
cualquier k.}

% ---------------------------------------------------------------------
\begin{frame}[fragile]{Rebanado con paso y rebanado como destino}
\begin{lstlisting}
v = [10, 20, 30, 40, 50, 60]

v[1:4]      # [20, 30, 40]
v[::2]      # [10, 30, 50]     uno de cada dos
v[-2:]      # [50, 60]         los dos últimos
v[::-1]     # [60, 50, 40, 30, 20, 10]
v[2:100]    # [30, 40, 50, 60] no falla: recorta al tamaño real

v[1:3] = [0, 0]     # el rebanado también sirve del lado izquierdo
del v[::2]          # y se puede eliminar por rebanado
\end{lstlisting}

Un rebanado de lista \textbf{siempre produce una lista nueva}. Recuérdelo: en
NumPy no será así.
\end{frame}
\note{v[:] es el modismo tradicional para copiar una lista completa; hoy se
prefiere v.copy() por claridad. Ambos siguen siendo copias superficiales.}

% ---------------------------------------------------------------------
\begin{frame}[fragile]{Recorrer una lista}
\begin{lstlisting}
medidas = [20.1, 20.4, 19.8, 20.0]

for x in medidas:                    # lo idiomático: el elemento
    print(x)

for i, x in enumerate(medidas):      # cuando hace falta el índice
    print(f"medida {i}: {x} °C")

for t, y in zip(tiempos, alturas):   # dos secuencias en paralelo
    print(t, y)

for i in range(len(medidas)):        # correcto, pero rara vez necesario
    print(medidas[i])
\end{lstlisting}
\end{frame}
\note{El último patrón es el que traen quienes vienen de C o Fortran.
Funciona, pero en Python se prefiere iterar sobre el objeto. enumerate y zip
son las dos herramientas que cubren casi todos los casos restantes.}

% ---------------------------------------------------------------------
\begin{frame}[fragile]{Comprensiones de lista}
Construir una lista a partir de otra, en una sola expresión.

\begin{lstlisting}
t = [0.0, 0.5, 1.0, 1.5, 2.0]

# y = ½ g t²  para cada instante
y = [0.5 * 9.81 * ti**2 for ti in t]

# con filtro
positivos = [x for x in medidas if x > 20.0]

# equivalente explícito de la primera
y = []
for ti in t:
    y.append(0.5 * 9.81 * ti**2)
\end{lstlisting}

Son más breves y algo más rápidas que el bucle, pero \textbf{no} dejan de ser
bucles del intérprete.
\end{frame}
\note{Este matiz importa para lo que sigue: la comprensión de lista no es
vectorización. En la medición que veremos, la comprensión y el bucle explícito
tardan prácticamente lo mismo.}

% ---------------------------------------------------------------------
\begin{frame}[fragile]{Listas anidadas: matrices improvisadas}
\begin{lstlisting}
M = [[1, 2, 3],
     [4, 5, 6],
     [7, 8, 9]]

M[1][2]        # 6      fila 1, columna 2
len(M)         # 3      número de filas
len(M[0])      # 3      número de columnas
[fila[0] for fila in M]    # [1, 4, 7]  primera columna, a mano
\end{lstlisting}

Con listas, una matriz es una lista de listas: las filas son objetos
independientes y \textbf{no hay garantía de que tengan la misma longitud}.
\end{frame}
\note{Pedirles que noten lo incómodo que resulta extraer una columna. En NumPy
será M[:, 0]. Esa incomodidad es el mejor argumento para introducir el
ndarray.}

% ---------------------------------------------------------------------
\begin{frame}[fragile]{Una trampa clásica}
\begin{lstlisting}
mal = [[0] * 3] * 3       # ✗
mal[0][0] = 1
# [[1, 0, 0], [1, 0, 0], [1, 0, 0]]   ¡las tres filas son el mismo objeto!

bien = [[0] * 3 for _ in range(3)]    # ✓
bien[0][0] = 1
# [[1, 0, 0], [0, 0, 0], [0, 0, 0]]
\end{lstlisting}

\code{* 3} repite \textbf{la referencia}, no el contenido. Es el mismo
asunto del alias, en otra forma.
\end{frame}
\note{Este error aparece siempre que alguien inicializa una matriz de ceros con
listas. En NumPy se resuelve con np.zeros((3,3)), que además reserva la
memoria de una vez.}

% ---------------------------------------------------------------------
\begin{frame}[fragile]{Ejercicio guiado: estadísticos de una serie}
\begin{lstlisting}
medidas = [20.1, 20.4, 19.8, 25.6, 20.0, 19.9]   # °C

n     = len(medidas)
media = sum(medidas) / n                          # 20.967

var   = sum((x - media)**2 for x in medidas) / (n - 1)
desv  = var ** 0.5                                # 2.279

mayor = max(medidas)                              # 25.6
pos   = medidas.index(mayor)                      # 3
\end{lstlisting}

Todo esto funciona y es correcto. La pregunta es qué pasa cuando
\code{medidas} tiene diez millones de elementos y hay que repetirlo mil
veces.
\end{frame}
\note{Dejar planteada la pregunta y no responderla todavía. La respuesta es la
segunda mitad de la sesión. El valor 25.6 es un dato atípico deliberado: lo
vamos a filtrar más adelante con una máscara de NumPy.}

% ---------------------------------------------------------------------
\begin{frame}[fragile]{Ejercicio guiado: producto punto}
\begin{lstlisting}
u = [1.0, 2.0, 3.0]
w = [4.0, 5.0, 6.0]

# versión explícita
punto = 0.0
for i in range(len(u)):
    punto += u[i] * w[i]

# versión idiomática
punto = sum(ui * wi for ui, wi in zip(u, w))      # 32.0
\end{lstlisting}

Funciona para vectores de cualquier dimensión, pero el bucle lo ejecuta el
intérprete, elemento por elemento.
\end{frame}
\note{Buen momento para preguntar cómo escribirían el producto matriz-vector
con listas. Tres bucles anidados y ninguna claridad. Ese es el punto de
quiebre.}

% =====================================================================
\seccion{2. NumPy y el ndarray}{El arreglo homogéneo de la computación científica}

% ---------------------------------------------------------------------
\begin{frame}[fragile]{Qué es NumPy}
\textbf{NumPy} (\textit{Numerical Python}) es la biblioteca sobre la que
descansa prácticamente todo el cómputo numérico en Python: SciPy, Matplotlib,
pandas, scikit-learn, Astropy y las bibliotecas de análisis de datos
experimentales.

\smallskip
Aporta un tipo de dato: el \textbf{ndarray}, un arreglo
\textbf{n-dimensional, homogéneo y de tamaño fijo}.

\begin{lstlisting}
import numpy as np        # la convención universal; úsela siempre así

a = np.array([3.2, 1.7, 4.8, 0.5])
\end{lstlisting}

Instalación: \code{pip install numpy} o \code{conda install numpy}.
\end{frame}
\note{Mencionar que NumPy nació en 2005 de la fusión de Numeric y numarray,
obra de Travis Oliphant, y que hoy es un proyecto con financiación
institucional. Para estudiantes de física, subrayar que Astropy y todo el
análisis de datos de LIGO están construidos sobre ndarray.}

% ---------------------------------------------------------------------
\begin{frame}{La diferencia está en la memoria}
\figura[0.70]{memoria-lista-vs-array}
\end{frame}
\note{Este es el diagrama central de la sesión. Si entienden esta diapositiva,
el resto de NumPy son detalles de sintaxis: la homogeneidad y la contigüidad
explican la velocidad, el ahorro de memoria, la existencia de dtype y la
rigidez del tamaño fijo.}

% ---------------------------------------------------------------------
\begin{frame}{Cuánto se gana, en números}
\figura[0.62]{rendimiento}
\end{frame}
\note{Los datos son de una medición propia con N = 1 000 000. Conviene
repetirla en vivo con \%timeit en Jupyter delante de ellos: el número exacto
cambia con la máquina, la relación no. Explicar que la ganancia no viene de
que NumPy "sea C", sino de que el bucle completo baja a código compilado una
sola vez.}

% ---------------------------------------------------------------------
\begin{frame}[fragile]{Crear arreglos}
\begin{lstlisting}
np.array([1.0, 2.0, 3.0])       # desde una lista
np.array([[1, 2], [3, 4]])      # desde listas anidadas: 2D

np.zeros((3, 4))                # matriz 3x4 de ceros
np.ones(5)                      # [1. 1. 1. 1. 1.]
np.full((2, 2), 9.81)           # relleno con un valor
np.eye(3)                       # identidad 3x3
np.empty(4)                     # sin inicializar: contiene basura

np.arange(0, 1, 0.25)           # [0.   0.25 0.5  0.75]   fin excluido
np.linspace(0, 1, 5)            # [0.   0.25 0.5  0.75 1. ]  fin incluido
\end{lstlisting}

Para muestrear un intervalo físico use \textbf{\code{linspace}}: se
especifica cuántos puntos, no el paso, y no acumula error de redondeo.
\end{frame}
\note{arange con paso flotante es una fuente conocida de sorpresas: por
acumulación de redondeo puede incluir o excluir el último punto de manera
impredecible. linspace calcula cada punto a partir de los extremos y no tiene
ese problema.}

% ---------------------------------------------------------------------
\begin{frame}[fragile]{Atributos de un arreglo}
\begin{lstlisting}
M = np.array([[1.0, 2.0, 3.0],
              [4.0, 5.0, 6.0]])

M.shape       # (2, 3)     tamaño de cada eje
M.ndim        # 2          número de ejes
M.size        # 6          total de elementos
M.dtype       # float64    tipo de TODOS los elementos
M.itemsize    # 8          bytes por elemento
M.nbytes      # 48         bytes del bloque de datos
\end{lstlisting}

El \code{dtype} se fija al crear el arreglo. Se puede pedir
explícitamente:

\begin{lstlisting}
np.array([1, 2, 3], dtype=np.float64)
np.zeros(1000, dtype=np.int32)
\end{lstlisting}
\end{frame}
\note{Advertir sobre la asignación silenciosa: si a es de dtype int y se le
asigna a[0] = 3.7, NumPy trunca a 3 sin avisar. En trabajo científico conviene
crear los arreglos como float64 desde el principio.}

% ---------------------------------------------------------------------
\begin{frame}[fragile]{Aritmética vectorizada}
Los operadores actúan \textbf{elemento a elemento}, sin bucle explícito.

\begin{lstlisting}
a = np.array([1.0, 2.0, 3.0])
b = np.array([4.0, 5.0, 6.0])

a + b         # [5. 7. 9.]
a * b         # [ 4. 10. 18.]    producto elemento a elemento
a ** 2        # [1. 4. 9.]
2 * a + 1     # [3. 5. 7.]
a @ b         # 32.0             producto punto (matmul)

np.sqrt(a)    # ufuncs: sqrt, sin, cos, exp, log, abs...
\end{lstlisting}

\begin{cita}
Una fórmula de física se escribe casi igual que en el papel, y se aplica a todo
el conjunto de datos de una vez.
\end{cita}
\end{frame}
\note{Contrastar explícitamente con la lista: [1,2] * 2 repite,
np.array([1,2]) * 2 duplica los valores. Y recordar que a * b NO es el
producto matricial: para eso está @ o np.matmul.}

% ---------------------------------------------------------------------
\begin{frame}{Broadcasting}
\figura[0.68]{broadcasting}
\end{frame}
\note{El caso práctico más frecuente en el curso: restar la media de cada
columna a una matriz de datos, M - M.mean(axis=0). Formas (n,k) y (k,) se
combinan sin escribir un solo bucle. Vale la pena hacerlo en vivo.}

% ---------------------------------------------------------------------
\begin{frame}[fragile]{Indexación y rebanado}
La sintaxis de listas se conserva, y se extiende a varios ejes con comas.

\begin{lstlisting}
a = np.arange(10)          # [0 1 2 ... 9]
a[3]                       # 3
a[2:5]                     # [2 3 4]
a[::-1]                    # invertido

M = np.arange(1, 10).reshape(3, 3)

M[1, 2]        # 6      fila 1, columna 2   (no M[1][2])
M[0]           # [1 2 3]     primera fila
M[:, 0]        # [1 4 7]     primera COLUMNA, en una expresión
M[0:2, 1:3]    # submatriz
\end{lstlisting}

Extraer una columna, que con listas exigía una comprensión, aquí es
\code{M[:, 0]}.
\end{frame}
\note{M[1][2] funciona pero crea un arreglo intermedio; M[1, 2] es el modismo
correcto y más eficiente. Con arreglos de más dimensiones la diferencia se
vuelve importante.}

% ---------------------------------------------------------------------
\begin{frame}{Rebanar un arreglo devuelve una vista}
\figura[0.66]{vista-vs-copia}
\end{frame}
\note{Este es el segundo punto donde se pierden. Aclarar que la vista es una
característica deseada: permite trabajar por bloques sobre arreglos enormes
sin duplicar memoria. Y que .copy() es explícito justamente para que la copia
sea una decisión, no un accidente.}

% ---------------------------------------------------------------------
\begin{frame}[fragile]{Máscaras booleanas}
Comparar un arreglo produce otro arreglo, de booleanos. Ese arreglo sirve para
seleccionar.

\begin{lstlisting}
medidas = np.array([20.1, 20.4, 19.8, 25.6, 20.0, 19.9])

medidas > 20.0            # [ True  True False  True False False]
medidas[medidas > 20.0]   # [20.1 20.4 25.6]

# descartar el dato atípico: dentro de 2 desviaciones estándar
mu, s = medidas.mean(), medidas.std(ddof=1)
buenas = medidas[np.abs(medidas - mu) < 2 * s]
buenas.mean()             # 20.04   (antes era 20.967)

np.sum(medidas > 20.0)    # 3       cuántas cumplen
np.where(medidas > 20.0)  # índices donde se cumple
\end{lstlisting}
\end{frame}
\note{Este es probablemente el modismo más útil de todo NumPy para un físico:
filtrar datos experimentales por criterio, sin bucles ni condicionales.
Comparar con lo que costaría hacerlo con listas y un if dentro de un for.}

% ---------------------------------------------------------------------
\begin{frame}{Reducciones y el parámetro axis}
\figura[0.66]{ejes-2d}
\end{frame}
\note{Reducciones disponibles: sum, mean, std, var, min, max, argmin, argmax,
prod, any, all, cumsum. Todas aceptan axis. argmin y argmax devuelven el
índice, no el valor: son las que se usan para localizar un pico.}

% ---------------------------------------------------------------------
\begin{frame}[fragile]{Forma y reorganización}
\begin{lstlisting}
a = np.arange(12)

M = a.reshape(3, 4)      # 3 filas, 4 columnas — es una VISTA de a
M.T                      # transpuesta, (4, 3) — también una vista
a.reshape(-1, 2)         # el -1 significa "calcúlelo usted"
M.ravel()                # de vuelta a 1D

np.concatenate([a, a])           # unir arreglos
np.vstack([f1, f2])              # apilar como filas
np.hstack([c1, c2])              # apilar como columnas
\end{lstlisting}

Un ndarray tiene \textbf{tamaño fijo}: no existe \code{append} barato. Si
necesita crecer elemento a elemento, \textbf{construya con una lista y
convierta al final}.
\end{frame}
\note{np.append existe pero reserva y copia todo el arreglo en cada llamada:
usarlo dentro de un bucle es cuadrático. El patrón correcto es acumular en una
lista de Python y llamar np.array una sola vez, o preasignar con np.zeros si
se conoce el tamaño.}

% ---------------------------------------------------------------------
\begin{frame}[fragile]{Aplicación: tiro parabólico vectorizado}
\begin{lstlisting}
import numpy as np

g, v0 = 9.81, 20.0
theta = np.radians(45.0)

T = 2 * v0 * np.sin(theta) / g          # tiempo de vuelo: 2.883 s
t = np.linspace(0.0, T, 5)              # [0. 0.721 1.442 2.162 2.883]

x = v0 * np.cos(theta) * t              # [0. 10.19 20.39 30.58 40.77]
y = v0 * np.sin(theta) * t - 0.5*g*t**2 # [0.  7.65 10.19  7.65  0.  ]

print(f"alcance {x[-1]:.2f} m, altura máxima {y.max():.2f} m")
\end{lstlisting}

Cinco puntos o cinco millones: el código es el mismo. Solo cambia el argumento
de \code{linspace}.
\end{frame}
\note{Con t = np.linspace(0, T, 500) y matplotlib se obtiene la trayectoria
completa en dos líneas más. Es un buen cierre en vivo: plt.plot(x, y);
plt.show()}

% ---------------------------------------------------------------------
\begin{frame}[fragile]{Aplicación: integrar datos discretos}
\begin{lstlisting}
x = np.linspace(0.0, 1.0, 5)
F = 3 * x**2                      # fuerza en función de la posición

W = np.trapezoid(F, x)            # 1.03125   (NumPy >= 2.0)
\end{lstlisting}

El valor exacto de $\int_0^1 3x^2\,dx$ es 1. La diferencia no es un error de
programación: es el \textbf{error de discretización} del método del trapecio
con cinco puntos.

\smallskip
Con \code{np.linspace(0, 1, 1001)} el resultado cae a 1{,}0000005.
\end{frame}
\note{Punto conceptual importante para físicos: el computador no integra,
aproxima. Aumentar el número de puntos reduce el error de truncamiento pero
aumenta el de redondeo acumulado. En versiones anteriores a NumPy 2.0 la
función se llama np.trapz.}

% =====================================================================
\seccion{3. Comparación y criterio}{}

% ---------------------------------------------------------------------
\begin{frame}[fragile]{Lista frente a ndarray}
\begin{center}
\rowcolors{2}{fondoClaro}{white}
\small
\begin{tabularx}{0.95\textwidth}{lXX}
\rowcolor{borgonaOscuro}
\color{white}\textbf{Aspecto} & \multicolumn{1}{l}{\color{white}\textbf{\texttt{list}}} & \multicolumn{1}{l}{\color{white}\textbf{\texttt{ndarray}}} \\
Tipo de elementos    & heterogéneo               & homogéneo (\texttt{dtype} único) \\
Memoria              & referencias dispersas     & bloque contiguo \\
Tamaño               & variable, \texttt{append} barato & fijo, crecer es caro \\
\texttt{a + b}       & concatena                 & suma elemento a elemento \\
\texttt{a * 2}       & repite                    & duplica los valores \\
Rebanado             & devuelve una copia        & devuelve una \textbf{vista} \\
Varias dimensiones   & listas anidadas, a mano   & nativo, con \texttt{shape} \\
Operar sin bucle     & no                        & sí, vectorizado \\
Velocidad numérica   & referencia                & $\approx 50\times$ más rápido \\
Dependencias         & ninguna                   & requiere NumPy \\
\end{tabularx}
\end{center}
\end{frame}
\note{La fila de la vista y la del rebanado son las dos que más confunden.
Conviene volver sobre ellas al final del taller práctico.}

% ---------------------------------------------------------------------
\begin{frame}{Cuándo usar cada uno}
\small
\subt{Lista}
\begin{itemize}
  \item Datos heterogéneos, o de tipo no numérico.
  \item Colecciones que \textbf{crecen} de tamaño desconocido: leer un archivo línea por línea, acumular resultados.
  \item Pocos elementos, donde la claridad importa más que la velocidad.
\end{itemize}

\subt{ndarray}
\begin{itemize}
  \item Cualquier cálculo numérico sobre muchos datos.
  \item Álgebra lineal, series temporales, mallas, imágenes, campos.
  \item Cuando se va a usar SciPy, Matplotlib, pandas o Astropy: todas esperan arreglos.
\end{itemize}
\begin{cita}
Patrón habitual: se acumula en una lista mientras se leen los datos, y se
convierte a arreglo cuando empieza el cálculo.
\end{cita}
\end{frame}
\note{Ese patrón final es el que van a usar en casi todos los talleres del
curso: leer un archivo de laboratorio a listas, np.array al terminar la
lectura, y de ahí en adelante todo vectorizado.}

% ---------------------------------------------------------------------
\begin{frame}[fragile]{Errores frecuentes}
\begin{enumerate}\small
  \item \code{v = v.sort()} — deja \code{v} en \code{None}.
  \item Confundir alias con copia, y modificar sin querer el original.
  \item Esperar que \code{lista1 + lista2} sume; concatena.
  \item Esperar que \code{arreglo1 + arreglo2} concatene; suma.
  \item Modificar una vista de NumPy creyendo que era una copia.
  \item \code{[[0]*n]*m} para crear una matriz: las filas quedan compartidas.
  \item Usar \code{np.append} dentro de un bucle: copia todo en cada iteración.
  \item Escribir un bucle \code{for} sobre un ndarray: anula la ventaja de NumPy.
  \item Mezclar \code{int} y \code{float} en un arreglo y perder decimales por truncamiento.
\end{enumerate}
\end{frame}
\note{El punto 8 merece énfasis: iterar un ndarray elemento por elemento es
más lento que iterar una lista. Si aparece un for sobre un arreglo, casi
siempre existe una versión vectorizada.}

% ---------------------------------------------------------------------
\begin{frame}[fragile]{Ejercicios propuestos}
\subt{Con listas, sin NumPy}
\begin{enumerate}
  \item Dada una lista de temperaturas en $^\circ$C, construya otra en K, sin modificar la original.
  \item Escriba una función que reciba una lista y devuelva la media, la mediana y la desviación estándar muestral.
  \item Sin usar \code{reverse()} ni \code{[::-1]}, invierta una lista en el sitio intercambiando extremos.
  \item Dada una lista de listas que representa una matriz, escriba una función que devuelva su transpuesta.
  \item Verifique con \code{id()} la diferencia entre \code{b = a}, \code{b = a[:]} y \code{copy.deepcopy(a)}.
\end{enumerate}
\end{frame}
\note{El ejercicio 3 obliga a razonar sobre índices y sobre la mutación en el
sitio. El 4 es el que mejor motiva la transición a NumPy: comparen su solución
con M.T una vez visto NumPy.}

% ---------------------------------------------------------------------
\begin{frame}[fragile]{Ejercicios propuestos}
\subt{Con NumPy}
\begin{enumerate}\small
  \setcounter{enumi}{5}
  \item Repita los ejercicios 1, 2 y 4 con arreglos y compare el número de líneas.
  \item Genere 1000 puntos entre 0 y $2\pi$ y evalúe $f(x)=\sin(x)\,e^{-x/4}$. Halle el máximo y la posición donde ocurre (\code{argmax}).
  \item Dada una matriz de datos de $n$ mediciones por $k$ sensores, réstele la media de cada columna en una sola línea.
  \item Con una máscara booleana, descarte de una serie los valores que se alejen más de $3\sigma$ de la media, e informe cuántos descartó.
  \item Mida con \code{\%timeit} el producto punto de dos vectores de $10^6$ elementos, con lista y con arreglo. Reporte el factor.
  \item Verifique experimentalmente que \code{b = a[2:5]} es una vista y \code{b = a[2:5].copy()} no lo es.
\end{enumerate}
\end{frame}
\note{El 8 es la aplicación directa del broadcasting. El 10 y el 11 son los
que consolidan los dos conceptos centrales de la sesión: vectorización y
vistas. Sugerir que entreguen el 10 con la medición de su propia máquina.}

% ---------------------------------------------------------------------
\begin{frame}[fragile]{Lo que sigue}
\begin{itemize}
  \item \textbf{Matplotlib}: graficar los arreglos que acabamos de construir.
  \item \textbf{Lectura de datos}: \code{np.loadtxt}, \code{np.genfromtxt} y archivos de laboratorio.
  \item \textbf{SciPy}: integración, optimización, ajuste de curvas, ecuaciones diferenciales.
  \item \textbf{Álgebra lineal}: \code{np.linalg} — sistemas, autovalores, descomposiciones.
  \item \textbf{Números aleatorios}: \code{np.random} y simulación Monte Carlo.
\end{itemize}

\bigskip
Referencia recomendada: \textcolor{borgonaClaro}{\texttt{numpy.org/doc/stable}}
— especialmente \textit{NumPy: the absolute basics for beginners}.
\end{frame}
\note{Anunciar la práctica: leer un archivo de datos de laboratorio, limpiarlo
con máscaras, ajustar una recta y graficar. Reúne todo lo de hoy.}

% =====================================================================
\seccion{Gracias}{Preguntas y discusión}

\end{document}
